\section{Tensors of Specific Shapes}

  % notation used only in this file
  \providecommand{\llbracket}{[\![}
  \providecommand{\rrbracket}{]\!]}
  \providecommand{\brank}{\underline{\mathrm{rank}}}

  Even though there is no general algorithm for computing the rank, the rank is completely understood for tensors of certain shapes. The tables in the review (maximum ranks of $I \times J \times 2$ and $3 \times 3 \times 3$ tensors, typical ranks of $2 \times 2 \times 2$, $3 \times 3 \times 2$, $I \times J \times 2$, tall tensors, tensors with symmetric slices, etc.) are summaries of a handful of papers, mostly from the psychometrics literature. In this section we go through them, with the main statement and an outline of the proof. We will use the \hyperref[lem:slice-span]{slice-span characterization} and the \hyperref[thm:pencil]{rank of a regular pencil} throughout.

\subsection{Tensors with Two Slices: $2 \times 2 \times 2$}

  The first result of this kind is due to Kruskal, who noticed that the set of $2 \times 2 \times 2$ tensors of rank less than the maximum rank has positive volume (this is what makes the typical rank not unique). The proof that was written down by Ten Berge in 1991 is the following.\footnote{J. M. F. Ten Berge, \textit{Kruskal's polynomial for $2 \times 2 \times 2$ arrays and a generalization to $2 \times n \times n$ arrays}, Psychometrika 56 (1991). The exposition of Bremner and Hu, \textit{On Kruskal's theorem that every $3 \times 3 \times 3$ array has rank at most 5}, also contains complete proofs of these facts.} For a tensor $\mathcal{X} \in \mathbb{R}^{2 \times 2 \times 2}$ with frontal slices $X_1, X_2$, define
  \begin{equation}
    \Delta(\mathcal{X}) \coloneqq \text{discriminant of the quadratic } \lambda \mapsto \det(X_2 - \lambda X_1)
  \end{equation}
  which is a homogeneous polynomial of degree $4$ in the entries, known as Cayley's \textbf{hyperdeterminant}, and it is also called Kruskal's polynomial.

  \begin{theorem}[Kruskal; Ten Berge]
    \label{thm:222}
    Let $\mathcal{X} \in \mathbb{R}^{2 \times 2 \times 2}$ be nonzero.
    \begin{enumerate}
      \item $\mathrm{rank}(\mathcal{X}) \leq 3$.
      \item $\mathrm{rank}(\mathcal{X}) = 1$ if and only if all six of its slices (in all three directions) are singular, unless $\mathcal{X}$ is \textbf{superdiagonal}, i.e. has exactly two nonzero entries that never lie in the same slice, which has rank $2$.
      \item Suppose that some slice is nonsingular, and by permuting the directions and slices say that it is $X_1$. If $X_2$ is a scalar multiple of $X_1$, then the rank is $2$. Otherwise,
      \begin{equation}
        \mathrm{rank}(\mathcal{X}) = \begin{cases} 2 & \text{if } \Delta(\mathcal{X}) > 0 \\ 3 & \text{if } \Delta(\mathcal{X}) \leq 0 \end{cases}
      \end{equation}
      Over $\mathbb{C}$ the condition becomes $\Delta(\mathcal{X}) \neq 0$ for rank $2$.
    \end{enumerate}
    In particular both $2$ and $3$ are typical ranks over $\mathbb{R}$, and the generic rank over $\mathbb{C}$ is $2$.
  \end{theorem}
  \begin{proof}
    By Lemma \ref{lem:slice-span}, $\mathrm{rank}(\mathcal{X})$ is the smallest $n$ such that $X_1 = A D B^T$ and $X_2 = A E B^T$ with $A, B \in \mathbb{R}^{2 \times n}$ and $D, E$ diagonal $n \times n$.

    \textit{Part 2.} If the rank is one, then $X_1 = d \, a b^T$ and $X_2 = e \, a b^T$, so both slices (and by symmetry all six) are singular. Conversely, if all slices are singular and some slice is zero, then the other slice is $a b^T$ and $\mathcal{X} = a \circ b \circ e_1$. If no slice is zero, write $X_1 = a b^T$ with $b = (1, \lambda)^T$, so that $X_1 = [a \,|\, \lambda a]$ in terms of columns, and then singularity of the second column slice and of the horizontal slices force the columns of $X_2$ to be $\mu a$ and $\lambda \mu a$, giving $\mathcal{X} = a \circ (1, \lambda)^T \circ (1, \mu)^T$ with rank one. The only exception is when the argument ends with $X_1 = [a \,|\, 0]$ and $X_2 = [0 \,|\, d]$, which is superdiagonal, and it has rank $2$ since $X_1, X_2$ are not scalar multiples of the same rank-one matrix.

    \textit{Part 1.} If every slice is singular then the rank is at most $2$ by Part 2. Otherwise, a nonsingular slice can be taken to be $X_1$. Let $Y = X_2 X_1^{-1} = (y_{ij})$. By multiplying on the right by $X_1^{-1}$ (a change of basis), it suffices to decompose the pair $(I, Y)$. Define
    \begin{equation}
      A = \begin{bmatrix} 1 & 0 & y_{12} \\ 0 & 1 & y_{21} \end{bmatrix}, \quad
      D = \mathrm{diag}(1, 1, 0), \quad
      E = \mathrm{diag}(y_{11} - y_{12}, \; y_{22} - y_{21}, \; 1), \quad
      B^T = \begin{bmatrix} 1 & 0 \\ 0 & 1 \\ 1 & 1 \end{bmatrix} X_1
    \end{equation}
    Then a direct computation gives $A D B^T = X_1$ and $A E B^T = Y X_1 = X_2$, so there is a decomposition with three terms.

    \textit{Part 3.} Suppose $X_1$ is nonsingular and $X_2$ is not a scalar multiple of $X_1$. By Part 1 the rank is $2$ or $3$, and by Theorem \ref{thm:pencil}, it is $2$ if and only if $X_1^{-1} X_2$ is diagonalizable over the field, i.e. (since $X_2$ is not a multiple of $X_1$) if and only if $\det(X_2 - \lambda X_1)$ has two distinct roots in the field. Concretely, if $X_1 = d_1 a_1 b_1^T + d_2 a_2 b_2^T$ and $X_2 = e_1 a_1 b_1^T + e_2 a_2 b_2^T$, then $X_2 - (e_j / d_j) X_1$ is singular for $j = 1, 2$, and these two values of $\lambda$ are distinct. Conversely, if $\lambda_1 \neq \lambda_2$ are roots, then the matrices $X_2 - \lambda_j X_1$ are singular and nonzero, hence rank one, and we can write
    \begin{equation}
      u_1 v_1^T = \frac{X_2 - \lambda_2 X_1}{\lambda_1 - \lambda_2}, \qquad u_2 v_2^T = -\frac{X_2 - \lambda_1 X_1}{\lambda_1 - \lambda_2}
    \end{equation}
    which satisfy $X_1 = u_1 v_1^T + u_2 v_2^T$ and $X_2 = \lambda_1 u_1 v_1^T + \lambda_2 u_2 v_2^T$. Two distinct real roots exist if and only if the discriminant $\Delta(\mathcal{X})$ is positive, and two distinct complex roots if and only if $\Delta(\mathcal{X}) \neq 0$.

    Finally, $\{\Delta > 0\}$ and $\{\Delta < 0\}$ are both nonempty open sets (e.g. $X_1 = I$ and $X_2 = \mathrm{diag}(1, 2)$ versus $X_2 = \begin{bmatrix} 0 & 1 \\ -1 & 0 \end{bmatrix}$), so ranks $2$ and $3$ both occur with positive probability.
  \end{proof}

  When $\Delta(\mathcal{X}) = 0$ and $X_2$ is not proportional to $X_1$, the rank is $3$ even though such a tensor is a limit of rank-two tensors. For example $X_1 = I$ and $X_2 = \begin{bmatrix} 0 & 1 \\ 0 & 0 \end{bmatrix}$ has rank $3$ (since the Jordan block is not diagonalizable), but it is the limit as $a \to 0$ of $(I, \begin{bmatrix} 0 & 1 \\ a^2 & 0 \end{bmatrix})$ which has $\Delta = 4a^2 \neq 0$ and hence rank $2$. This is again the example of border rank two but rank three. Monte Carlo experiments with Gaussian entries give about $79\%$ rank two and $21\%$ rank three.

\subsection{Tensors with Two Slices: $I \times I \times 2$ and $I \times J \times 2$}

  The same idea generalizes to all pencils, and this is what is behind the $K = 2$ entries in the review's tables. The first result is about the maximum rank.

  \begin{theorem}[Kruskal; J\'aJ\'a]
    \label{thm:max-pencil}
    The maximum rank of an $I \times J \times 2$ tensor is $\min\{I, J\} + \min\{I, J, \lfloor \max\{I, J\}/2 \rfloor\}$. In particular, the maximum rank of an $I \times I \times 2$ tensor is $\lfloor 3I/2 \rfloor$.
  \end{theorem}
  \begin{proof}
    We prove the upper bound for $I \times I \times 2$ tensors whose pencil is regular, i.e. such that after a change of basis in mode $3$ the first slice $X_1$ is invertible. The general case is handled in the same way with the Kronecker canonical form of the pencil. By a change of basis in modes $1, 2$, we may replace $(X_1, X_2)$ by $(I, C)$ where $C = X_1^{-1} X_2$ is in rational (or Jordan) canonical form, which is block diagonal with blocks $C_1, \ldots, C_m$ of sizes $s_1, \ldots, s_m$, each the companion matrix of a polynomial. Then $(I, C)$ is the direct sum of the pairs $(I_{s_i}, C_i)$, so the rank is at most the sum of the ranks.

    Each block satisfies $\mathrm{rank}(I_s, C_s) \leq s + 1$, and $= s$ if $s = 1$. To see this, write the companion matrix as $C_s = C' + D$, where $D$ is the rank-one matrix that agrees with $C_s$ on the last column and is zero elsewhere, and $C'$ is the companion matrix of a polynomial with $s$ distinct real roots (any monic polynomial is a companion matrix after changing the last column, so we are free to choose these roots). Then
    \begin{equation}
      (I_s, C_s) = (0, D) + (I_s, C')
    \end{equation}
    where $(0, D)$ is rank one, and $(I_s, C')$ has rank $s$ by Theorem \ref{thm:pencil} since $C'$ is diagonalizable.

    So the rank is at most $n + \#\{i : s_i \geq 2\} \leq n + \lfloor n/2 \rfloor$, since blocks of size $1$ are free and there are at most $\lfloor n/2 \rfloor$ blocks of size at least $2$.
  \end{proof}

  This explains $2 \times 2 \times 2 \mapsto 3$, $3 \times 3 \times 2 \mapsto 4$ and $4 \times 4 \times 2 \mapsto 6$. Next we turn to the typical rank.

  \begin{theorem}[Ten Berge and Kiers]
    \label{thm:typical-pencil}
    The typical ranks of $I \times J \times 2$ tensors over $\mathbb{R}$ are the following.
    \begin{enumerate}
      \item If $I = J$, then the typical ranks are $\{I, I+1\}$.
      \item If $J < I < 2J$, then the typical rank is $I$.
      \item If $I \geq 2J$, then the typical rank is $2J$.
    \end{enumerate}
  \end{theorem}
  \begin{proof}
    \textit{Part 1.} With probability one, $X_1$ is invertible, so $\mathrm{rank}(\mathcal{X}) \geq \mathrm{rank}(X_1) = I$. Let $Y = X_1^{-1} X_2$, which has distinct eigenvalues with probability one, and hence is similar to the companion matrix of its characteristic polynomial.
    \begin{itemize}
      \item If all eigenvalues of $Y$ are real, which is an open condition and so has positive probability, then $Y$ is diagonalizable and the rank is $I$ by Theorem \ref{thm:pencil}.
      \item If $Y$ has a pair of complex eigenvalues, which is also an open condition, then $Y$ is not diagonalizable over $\mathbb{R}$, so the rank is more than $I$ by Theorem \ref{thm:pencil}.
      \item In both cases the rank is at most $I + 1$ by the companion matrix argument in the proof of Theorem \ref{thm:max-pencil}, applied to a single block.
    \end{itemize}
    Hence rank $I$ and rank $I+1$ both occur with positive probability, and no other value does.
  \end{proof}
  \begin{proof}
    \textit{Parts 2 and 3.} These are special cases of the tall tensor theorem proven below with $K = 2$.
  \end{proof}

  Ten Berge's empirical check of how large the set of rank-$n$ tensors is: for pairs of $n \times n$ Gaussian matrices, the fractions of $100$ samples for which $X_2 X_1^{-1}$ had $n$ distinct real eigenvalues were $76, 55, 22, 16, 2, 1, 0$ for $n = 2, \ldots, 8$. So rank $n$ is typical but it becomes increasingly rare. Since the typical ranks in these cases are not unique, we can see why the review lists the typical rank as $\{3, 4\}$ for $3 \times 3 \times 2$, and $\{I, I+1\}$ for $I \times I \times 2$.

  The same technique shows that the typical and maximum rank can be far apart.

  \begin{example}[$4 \times 4 \times 2$]
    The maximum rank of $4 \times 4 \times 2$ tensors is $6$, but the typical rank is $\{4, 5\}$. For instance, take $X_1 = I_4$ and
    \begin{equation}
      X_2 = \begin{bmatrix} 0 & -1 & 0 & 0 \\ 1 & 0 & 0 & 0 \\ 0 & 0 & 0 & -b \\ 0 & 0 & b & 0 \end{bmatrix}, \qquad b \neq 0
    \end{equation}
    whose eigenvalues $\pm i, \pm bi$ are complex, so the rank is at least $5$. For $b^2 \neq 1$, subtracting a suitable rank-one matrix $x y^T$ from $X_2$ gives a matrix with the real distinct eigenvalues $\pm 1, \pm b$, so the rank is exactly $5$. When $b = \pm 1$ this construction is not possible, and the rank is $6$ (shown by Thijsse). So the rank is $5$ with probability one when $b$ is random, but the maximal rank $6$ occurs on a set of measure zero.
  \end{example}

\subsection{$3 \times 3 \times 3$ Tensors}

  The case that is not a pencil but is still completely understood is $3 \times 3 \times 3$.

  \begin{theorem}[Kruskal]
    \label{thm:333}
    Every $3 \times 3 \times 3$ tensor over $\mathbb{R}$ or $\mathbb{C}$ has rank at most $5$. Combined with the fact that the generic rank of $3 \times 3 \times 3$ tensors over $\mathbb{C}$ is $5$, the maximum rank and typical rank are both $5$.
  \end{theorem}
  \begin{proof}
    This was stated by Kruskal in 1989 without proof, and a complete proof was written down by Bremner and Hu (2010), following a simplification by Rocci.\footnote{\url{http://www.cecm.sfu.ca/~jha107/BH333rank5v7.pdf}} Let $\mathcal{T} = [A \,|\, B \,|\, C]$ be the tensor with frontal slices $A, B, C$. The lower bound for the generic rank is Table 2 of Comon et al., discussed later; we outline the upper bound.

    \textit{Reduction to two singular slices.} If some slice is zero, then $\mathcal{T}$ is a $3 \times 3 \times 2$ tensor and has rank at most $4$ by Theorem \ref{thm:max-pencil}. Otherwise, if at least two slices are nonsingular, say $C$ is nonsingular, then $\det(A - \lambda C)$ is a cubic polynomial in $\lambda$ with leading coefficient $\pm \det C \neq 0$. Since the degree is odd, it has a real root, so replacing $A$ by $A - \lambda C$ (a change of basis in mode $3$) makes it singular. We do the same for $B$. So we may assume that $A$ and $B$ are singular.

    \textit{Reduction to slices of rank at least two.} If some slice has rank $\leq 1$, say $A$, then $\mathcal{T} = [A|0|0] + [0|B|C]$ where the first term is rank one and the second is a $3 \times 3 \times 2$ tensor of rank at most $4$. So we may assume $A, B$ have rank exactly $2$, and $C$ has rank $2$ or $3$.

    \textit{Case analysis.} We use two lemmas of Kruskal, whose proofs are elementary.
    \begin{itemize}
      \item \textbf{One-Edge Lemma.} If two slices $D, E$ have a common null vector, $Dx = Ex = 0$ (or $D^T x = E^T x = 0$), then $\mathrm{rank}(\mathcal{T}) \leq 5$. \textit{Proof:} change basis in mode $2$ so that $x = e_1$. Then the first column of $D$ and $E$ vanishes, so $\mathcal{T}$ is the sum of a tensor that is nonzero only in the first column of the third slice (rank one), and a tensor of size $3 \times 2 \times 3$, which has rank at most $4$ as a (transposed) $3 \times 3 \times 2$ tensor.
      \item \textbf{Two-Edge Lemma.} If two slices $D, E$ have $Dx = 0$, $y^T D = 0$ and $y^T E x \neq 0$, then $\mathrm{rank}(\mathcal{T}) \leq 5$. \textit{Proof:} choose bases with $x, y$ as first vectors so that $D$ has zero first row and column, and the $(1,1)$ entry $\alpha = y^T E x$ of $E$ is nonzero. Then the first row and column of $E$ (and, after replacing $C$ by $C + E$ if necessary, of $C$) can be matched by a rank-one matrix $\alpha^{-1} u v^T$. Subtracting these two rank-one tensors leaves a $2 \times 2 \times 2$ tensor, which has rank at most $3$ by Theorem \ref{thm:222}.
    \end{itemize}
    If $A, B, C$ all have rank $2$, take the null vectors $x_i$ and left null vectors $y_i$ of the three slices and form $X = [x_1 | x_2 | x_3]$, $Y = [y_1 | y_2 | y_3]$, so that $Y^T [A|B|C] X$ has a zero pattern. If the Two-Edge Lemma applies, we are done. If not, then all the diagonal entries $y_i^T M x_i$ vanish, and we split into cases on the ranks of $X$ and $Y$. When two columns of $X$ are dependent, we use the One-Edge Lemma; when $X$ has rank $2$ one finds that the first rows of the three slices are dependent and removes them to get a $2 \times 3 \times 3$ tensor of rank at most $4$; and when both $X, Y$ are invertible one subtracts three explicit rank-one tensors to be left with two more. When $C$ is invertible, a similar case analysis applies, now using the cases $\alpha A + \beta B + C$ singular or not. Each case exhibits a decomposition with at most five terms.
  \end{proof}

  Over a field that is not algebraically closed or $\mathbb{R}$, the statement can fail: over $\mathbb{F}_2$, there exist $3 \times 3 \times 3$ arrays of rank $6$ (von zur Gathen). The odd degree $3$ is used in the \textit{reduction to two singular slices} step, where over $\mathbb{R}$ we need the cubic $\det(A - \lambda C)$ to have a real root.

\subsection{Tall Tensors}

  The case where one dimension is large compared to the others was settled by Ten Berge. Let $\mathcal{X} \in \mathbb{R}^{I \times J \times K}$ with $I \geq J \geq K$. We call it
  \begin{itemize}
    \item \textbf{very tall} if $I \geq JK$,
    \item \textbf{tall} if $JK - J < I < JK$,
    \item \textbf{compact} if $I \leq JK - J$.
  \end{itemize}

  \begin{theorem}[Ten Berge, 2000]
    \label{thm:tall}
    Let $\mathcal{X} \in \mathbb{R}^{I \times J \times K}$ with $I \geq J \geq K$.
    \begin{enumerate}
      \item If $\mathcal{X}$ is very tall, then the typical rank (and maximum rank) is $JK$.
      \item If $\mathcal{X}$ is tall, then the typical rank is $I$.
      \item If $I = JK - J$, then the rank is $I$ or $I + 1$ almost surely, and rank $I$ occurs with positive probability. Moreover, if $K > 2$ and $J$ is odd then the typical rank is exactly $I$.
    \end{enumerate}
  \end{theorem}

  \begin{proof}
    \textit{Part 1.} Let $X_1, \ldots, X_K \in \mathbb{R}^{I \times J}$ be the frontal slices. Almost surely, the $I \times JK$ matrix $[X_1 | \cdots | X_K]$ has rank $JK$, so by the flattening bound $\mathrm{rank}(\mathcal{X}) \geq JK$. For the upper bound, write $X_k = A_k B_k^T$ with $A_k = X_k$ and $B_k = I_J$, and put $A = [A_1 | \cdots | A_K] \in \mathbb{R}^{I \times JK}$, $B = [B_1 | \cdots | B_K] \in \mathbb{R}^{J \times JK}$, $D_k = \mathrm{diag}(0, \ldots, 0, 1_J, 0, \ldots, 0)$ with the ones in the $k$th block. Then $X_k = A D_k B^T$ for all $k$, which is a decomposition with $JK$ terms.
  \end{proof}

  \begin{proof}
    \textit{Part 2.} The lower bound $I$ is again from the flattening $[X_1 | \cdots | X_K]$ having rank $I$ almost surely (since $I < JK$). For the upper bound, we need the following.

    \textit{Claim.} If $\mathrm{rank}(\mathcal{X}) = I$ and $JK > I$, then almost surely there is a nonsingular $S \in \mathbb{R}^{I \times I}$ such that $S X_1, \ldots, S X_K$ all have proportional rows. Indeed, if $X_k = A D_k B^T$ with $A \in \mathbb{R}^{I \times I}$, then $A$ must be invertible almost surely (otherwise the rows of $[X_1 | \cdots | X_K]$ would not span $\mathbb{R}^I$), and $S = A^{-1}$ gives $S X_k = D_k B^T$, whose $i$th row is the $i$th row of $B^T$ scaled by $(D_k)_{ii}$. Conversely, such an $S$ gives a rank $I$ decomposition $X_k = S^{-1} D_k (S X_1)$.

    So we need to construct $S$. Take the weights $d_{ik} = i^{k-1}$ (the Vandermonde matrix), and ask that each row $s_i$ of $S$ satisfies
    \begin{equation}
      s_i \big( i \, X_1 - X_2 \big) = s_i \big( i^2 X_1 - X_3 \big) = \cdots = s_i \big( i^{K-1} X_1 - X_K \big) = 0
    \end{equation}
    That is, $s_i$ is orthogonal to the columns of the $I \times (JK - J)$ matrix $W_i = [i X_1 - X_2 \,|\, i^2 X_1 - X_3 \,|\, \cdots \,|\, i^{K-1} X_1 - X_K]$. Since $I > JK - J$, such a nonzero $s_i$ always exists, e.g. $s_i = H_i e_i$ where $H_i = I - W_i (W_i^T W_i)^+ W_i^T$ is the projection onto the orthogonal complement of the columns of $W_i$. It remains to show that the resulting matrix $S$ is nonsingular almost surely. Its determinant is an analytic function of the entries of $\mathcal{X}$, so it is either identically zero or nonzero almost everywhere (Fisher's theorem). It is not identically zero since there are tensors for which the proportionality condition already holds with $S = I$, and for those the construction returns $S = I$. Hence $\mathrm{rank}(\mathcal{X}) \leq I$ almost surely.
  \end{proof}

  \begin{proof}
    \textit{Part 3.} Rank less than $I$ would make the $I$ rows of $[X_1 | \cdots | X_K]$ dependent, which has probability zero. Adding a random horizontal slice gives a $(JK - J + 1) \times J \times K$ tensor which is tall, so has rank $I + 1$ almost surely by (2), and removing a slice cannot increase the rank, so $\mathrm{rank}(\mathcal{X}) \leq I + 1$.

    To see that rank $I$ occurs with positive probability, transform the first $I$ columns of $[X_1 | \cdots | X_K]$ to $I_I$ by a change of basis in mode $1$, so that the last frontal slice becomes $[Y_{1K}; \ldots; Y_{K-1,K}]$ stacked from $K - 1$ blocks of size $J \times J$, while the other slices are block selectors. If all of the $Y_{jK}$ have real eigendecompositions $Y_{jK} = K_j \Lambda_j K_j^{-1}$ (an open condition), then they can be simultaneously diagonalized in the corresponding rows, which gives an explicit decomposition with $I = (K-1) J$ terms.

    Now suppose $K > 2$ and $J$ is odd. Then even without real eigendecompositions we can construct $S$ as before. For each $i$ we need $b_i$, $d_{i1}, \ldots, d_{iK}$ with
    \begin{equation}
      b_i \Big[ i \, Y_{1K} + i^2 Y_{2K} + \cdots + i^{K-1} Y_{K-1,K} - d_{iK} I_J \Big] = 0
    \end{equation}
    which says that $d_{iK}$ is an eigenvalue and $b_i$ a left eigenvector of the $J \times J$ matrix $\sum_k i^k Y_{kK}$. Since $J$ is odd, every real $J \times J$ matrix has a real eigenvalue, so we can solve this for each $i$ separately, and nonsingularity of the resulting $S$ holds almost surely by the same analyticity argument. Hence the typical rank is $I$.
  \end{proof}

  The hypotheses in (3) are sharp. For $K = 2$ the weights $i^{k-1}$ do not differ across blocks, so all the weighted combinations have the same eigenvectors and the argument breaks down, and indeed the typical ranks are $\{I, I+1\}$ by Theorem \ref{thm:typical-pencil}. For $J$ even it can happen that every linear combination of the $Y_{jK}$ has only complex eigenvalues, and then no real $b_i$ exists. For instance, for $8 \times 4 \times 3$ tensors, such arrays are rare but arise with positive probability, and both ranks $8$ and $9$ are typical. Ten Berge's table of results for $K = 2, 3$ is the following.
  \begin{center}
    \begin{tabular}{c|ccc@{\qquad}c|ccc}
      $K = 2$ & $J = 2$ & $J = 3$ & $J = 4$ & $K = 3$ & $J = 3$ & $J = 4$ & $J = 5$ \\ \hline
      $I = 2$ & $\{2, 3\}$ & $3$ & $4$ & $I = 5$ & ? & ? & ? \\
      $I = 3$ & $3$ & $\{3, 4\}$ & $4$ & $I = 6$ & $6$ & ? & ? \\
      $I = 4$ & $4$ & $4$ & $\{4, 5\}$ & $I = 7$ & $7$ & ? & ? \\
      $I = 5$ & $4$ & $5$ & $5$ & $I = 8$ & $8$ & $\{8, 9\}$ & ? \\
      $I = 6$ & $4$ & $6$ & $6$ & $I = 9$ & $9$ & $9$ & ? \\
    \end{tabular}
  \end{center}
  (Entries with $I < J$ follow from the case $I \geq J$ by permuting the modes.) The compact tensors with $I < JK - J$ are not covered by this method, and their typical ranks are only known from numerical computation, which we turn to next.

\subsection{Generic Ranks by Computing Dimensions}

  For the compact cases, the rank can be determined numerically, and the idea is the following. Consider the map
  \begin{equation}
    \varphi : \{ a_r, b_r, c_r \}_{r=1}^R \mapsto \sum_{r=1}^R a_r \circ b_r \circ c_r
  \end{equation}
  from a space of dimension $M = R(I + J + K - 2)$ (we subtract $2$ per term since scalings can be moved between the three factors) to $\mathbb{R}^{I \times J \times K}$.

  \begin{theorem}[Comon, Ten Berge, De Lathauwer, Castaing, 2009]
    Let $D(R)$ be the rank of the Jacobian of $\varphi$ at a random point. Then $D(R)$ is the dimension of the closure of the set of tensors of rank at most $R$. If $D(R) = IJK$, then $R$ is a typical rank, and the smallest such $R$ is the smallest typical rank (the generic rank over $\mathbb{C}$).
  \end{theorem}
  \begin{proof}
    The set of tensors of rank at most $R$ is the image of an irreducible variety under $\varphi$, so its closure is irreducible, and the dimension of an irreducible variety equals the dimension of its tangent space at a generic point. By Terracini's lemma, this tangent space is the image of the differential of $\varphi$, i.e. the column space of the Jacobian. If its rank is $IJK$, then the image contains an open set, which has positive volume. Increasing $R$ one at a time until the rank stops increasing gives the smallest $R$ for which the Jacobian has full rank. Explicitly, the rows of the Jacobian for the term $r$ are $I_I \otimes b_r \otimes c_r$, $a_r \otimes I_J \otimes c_r$, $a_r \otimes b_r \otimes I_K$.
  \end{proof}

  Running this for random parameters gives the following smallest typical ranks for cubical tensors of order $3$, which are the generic ranks over $\mathbb{C}$.
  \begin{center}
    \begin{tabular}{c|cccccccc}
      $N$ & $2$ & $3$ & $4$ & $5$ & $6$ & $7$ & $8$ & $9$ \\ \hline
      $N \times N \times N$ & $2$ & $5$ & $7$ & $10$ & $14$ & $19$ & $24$ & $30$ \\
    \end{tabular}
  \end{center}
  The naive dimension count gives the generic rank as at least $\lceil N^3 / (3N - 2) \rceil$, which equals $4$ for $N = 3$ and not $5$, so $3 \times 3 \times 3$ is \textit{defective}. The numerical results also include the tables for $2$-, $3$-, and $4$-slice tensors, for example the entries $\{2,3\}$ for $2 \times 2 \times 2$, $\{3, 4\}$ for $3 \times 3 \times 2$, $\{5, 6\}$ for $5 \times 3 \times 3$, which agree with all of the known results above. Moreover, the generic rank is a necessary condition for uniqueness: a tensor of rank larger than the generic value has infinitely many decompositions, whereas a tensor of smaller rank has almost surely finitely many.

\subsection{Tensors with Symmetric Slices}

  When fitting INDSCAL, the tensor has frontal slices that are symmetric, $\mathcal{X} \in \mathbb{R}^{I \times J \times J}$ with each $J \times J$ slice symmetric, and the factors are $A = B$. Symmetry often lowers the typical rank, and it is best understood by looking at what the slice-span lemma says.

  \begin{lemma}[Symmetric Slices]
    \label{prop:symmetric-slices}
    Let $\mathcal{X}$ have $K$ symmetric $J \times J$ frontal slices. Then $\mathrm{rank}(\mathcal{X}) \leq J(J+1)/2$. Moreover, if $K \geq J(J+1)/2$ then the typical rank is exactly $J(J+1)/2$.
  \end{lemma}
  \begin{proof}
    The space of symmetric $J \times J$ matrices has dimension $J(J+1)/2$ and is spanned by the rank-one matrices $e_i e_i^T$ and $(e_i + e_j)(e_i + e_j)^T$. The slices lie in this span, so by Lemma \ref{lem:slice-span}, $\mathrm{rank}(\mathcal{X}) \leq J(J+1)/2$. For the lower bound, when $K \geq J(J+1)/2$ generic symmetric slices span the whole space of symmetric matrices, so the flattening $X_{(3)}$ has rank $J(J+1)/2$.
  \end{proof}

  This explains the entries in the review where the symmetric typical rank is lower than the non-symmetric one: for $I \times 2 \times 2$ with $I \geq 4$, the symmetric typical rank is $3$ instead of $4$, and for $I \times 3 \times 3$ with $I \geq 6$ it is $6$ rather than $\min\{I, 9\}$. In particular, a $7 \times 3 \times 3$ array of symmetric slices is tall, but its typical rank is $6$ and not $7$, which means that the tall tensor theorem above does not apply to the symmetric case (symmetry does not arise with random sampling). When the slices are not too many, the typical ranks have been determined by Ten Berge, Sidiropoulos, and Rocci for arrays of order $I \times 2 \times 2$ and $I \times 3 \times 3$, and computed numerically by the Jacobian method of Comon et al. for others. For example, for $I \times I \times 2$ with symmetric slices the typical rank is again $\{I, I+1\}$, and for $3 \times 3 \times 3$ with symmetric slices it is $4$, compared to $5$ without symmetry.

  \begin{lemma}[Double Centering]
    The typical rank of $N_1 \times N_2 \times N_3$ tensors with all slices \textbf{double centered} (every row and column of every $N_2 \times N_3$ slice sums to zero) equals the typical rank of free $N_1 \times (N_2 - 1) \times (N_3 - 1)$ tensors.
  \end{lemma}
  \begin{proof}
    A double centered slice lies in $\mathbf{1}^\perp \otimes \mathbf{1}^\perp$, and a basis of $\mathbf{1}^\perp$ identifies this with $\mathbb{R}^{N_2 - 1} \otimes \mathbb{R}^{N_3 - 1}$. Applying the projection onto $\mathbf{1}^\perp$ in modes $2$ and $3$ fixes the tensor and maps rank-one terms to rank-one terms, so the rank can be computed using only centered factors. Therefore the rank is the same as that of the corresponding $N_1 \times (N_2 - 1) \times (N_3 - 1)$ tensor, and the change of basis is a linear isomorphism, which does not change the rank.
  \end{proof}

  Finally, for the fully symmetric case, the numerical method recovers the Alexander--Hirschowitz values: the generic symmetric rank of cubic forms in $N = 2, \ldots, 8$ variables is $2, 4, 5, 8, 10, 12, 15$, with the exceptional value $8$ for $N = 5$ (compared to the formula value $7$), and for quartics in $N = 2, \ldots, 8$ variables it is $3, 6, 10, 15, 21, 30, 42$.
